\documentclass[aspectratio=169]{beamer}

\usepackage[utf8]{inputenc}
\usepackage[T1]{fontenc}
\usepackage[spanish,es-tabla]{babel}
\usepackage{amsmath}
\usepackage{amssymb}
\usepackage{booktabs}
\usepackage{listings}
\usepackage{xcolor}

\usetheme{Madrid}
\usecolortheme{whale}
\setbeamertemplate{navigation symbols}{}
\setbeamertemplate{caption}[numbered]

\definecolor{codebg}{RGB}{245,245,245}
\lstset{
  basicstyle=\ttfamily\scriptsize,
  backgroundcolor=\color{codebg},
  breaklines=true,
  frame=single,
  columns=fullflexible,
  showstringspaces=false
}

\title[Envenenamiento de KG-RAG]{Ataques de envenenamiento de conocimiento a sistemas KG-RAG}
\subtitle{Reproducción del pipeline de ataque en caja negra de Zhao et al. (arXiv:2507.08862)}
\author{Ingrid Alicia Yábar Geldres}
\institute[UNI]{
  Universidad Nacional de Ingeniería (UNI)\\
  Facultad de Ingeniería Económica, Estadística y Ciencias Sociales\\
  Doctorado en Ciencias e Ingeniería Estadística\\
  \small ORCID: 0009-0003-9729-3259 \quad ingrid.yabar.g@uni.pe
}
\date{}

\begin{document}

\begin{frame}
  \titlepage
\end{frame}

\begin{frame}{Contenido}
  \tableofcontents
\end{frame}

% ============================================================
\section{Motivación}
% ============================================================

\begin{frame}{¿Por qué importa la seguridad de KG-RAG?}
  \begin{itemize}
    \item RAG permite que un LLM responda apoyándose en una fuente de conocimiento
      externa, en lugar de depender solo de lo memorizado en el entrenamiento.
    \item Cuando esa fuente es un \textbf{grafo de conocimiento (KG)}, el sistema
      puede seguir caminos de razonamiento explícitos y auditables entre
      entidades: KG-RAG.
    \item La misma propiedad que hace útil a un KG-RAG --- un grafo editable y
      actualizable --- es también su punto débil.
    \item Una sola tripleta falsa, bien conectada a la entidad correcta, puede
      desviar un camino de razonamiento multi-salto completo, sin alterar
      visiblemente el resto del grafo.
  \end{itemize}
\end{frame}

\begin{frame}{Ejemplo motivador}
  \begin{block}{Pregunta}
    ¿En qué país se filmó la película \emph{Manchester By The Sea}?
  \end{block}
  \begin{columns}
    \begin{column}{0.5\textwidth}
      \textbf{Antes del ataque}
      \begin{itemize}
        \item[] Manchester By The Sea \(\to\) \texttt{filmPlace} \(\to\) Manchester
        \item[] Manchester \(\to\) \texttt{locatedIn} \(\to\) Massachusetts
        \item[] Massachusetts \(\to\) \texttt{stateOf} \(\to\) United States
      \end{itemize}
    \end{column}
    \begin{column}{0.5\textwidth}
      \textbf{Después de insertar 1 arista falsa}
      \begin{itemize}
        \item[] Manchester \(\to\) \texttt{cityOf} \(\to\) England
        \item[] England \(\to\) \texttt{containedIn} \(\to\) United Kingdom
        \item[] $\Rightarrow$ el sistema puede responder \emph{United Kingdom}
      \end{itemize}
    \end{column}
  \end{columns}
  \vspace{0.5em}
  \small Objetivo de esta presentación: documentar una implementación propia de
  este ataque, de extremo a extremo, sobre un caso verificable a mano.
\end{frame}

% ============================================================
\section{Marco teórico}
% ============================================================

\begin{frame}{Grafo de conocimiento y caminos de relaciones}
  \begin{block}{Grafo de conocimiento}
    $G = (E, R, T)$, con $T \subseteq E \times R \times E$ el conjunto de
    tripletas $(e_h, r, e_t)$: una arista dirigida y etiquetada de $e_h$ a $e_t$.
  \end{block}
  \begin{block}{Camino de razonamiento y camino de relaciones}
    $p: e_0 \xrightarrow{r_1} e_1 \xrightarrow{r_2} \cdots \xrightarrow{r_l} e_l$
    \quad$\Rightarrow$\quad camino de relaciones asociado $w = (r_1, r_2, \dots, r_l)$
    (se descartan las entidades intermedias, se conserva solo la secuencia de relaciones).
  \end{block}
  \begin{block}{Anclaje (\emph{grounding})}
    Operación inversa: dado un camino de relaciones $w$ y una entidad de inicio,
    encontrar las entidades reales del grafo alcanzadas al recorrer $w$.
  \end{block}
\end{frame}

\begin{frame}{Pipeline de un sistema KG-RAG}
  \begin{block}{Etapa de recuperación}
    \[ G_q = \text{Retriever}(q; G) \tag{1} \]
    $G_q = (E_q, R_q, T_q)$: subgrafo extraído del grafo completo $G$ con solo
    la información relevante para responder la pregunta $q$.
  \end{block}
  \begin{block}{Etapa de generación}
    \[ a = f_\theta(\text{Prompt}[q, G_q]) \tag{2} \]
    $\theta$: parámetros del LLM generador; $\text{Prompt}[q, G_q]$: plantilla
    que combina la pregunta con el subgrafo recuperado; $a \in E$: respuesta
    final que produce el sistema.
  \end{block}
  \begin{alertblock}{Vulnerabilidad}
    Si $G$ contiene tripletas falsas bien conectadas, la recuperación puede
    incluirlas en $G_q$ sin distinguirlas de las legítimas.
  \end{alertblock}
\end{frame}

\begin{frame}{Modelo de amenaza (caja negra)}
  \[ \hat{A} = \text{KG-RAG}(q; \hat{G}), \qquad a^{*} \notin \hat{A}, \qquad
     \hat{G} = \{E, R, T \cup \hat{T}\} \tag{3} \]
  \begin{itemize}
    \item $\hat{T}$: tripletas adversarias insertadas; $\hat{G}$: grafo
      envenenado ($T \cup \hat{T}$); $\hat{A}$: respuestas que produce el
      sistema al operar sobre $\hat{G}$; $a^{*}$: respuesta correcta original.
    \item El atacante \textbf{no} conoce el recuperador, el LLM, ni la respuesta
      correcta $a^{*}$ de la pregunta atacada.
    \item Su única capacidad es \textbf{insertar} $\hat{T}$ en el grafo --- no
      puede borrar ni modificar tripletas existentes.
    \item Restricciones de sigilo: cada tripleta insertada reutiliza entidades y
      relaciones \emph{ya existentes} en el grafo, y el número de tripletas por
      respuesta adversaria está acotado por un presupuesto $K$.
  \end{itemize}
\end{frame}

\begin{frame}{Formalización: objetivo de extracción de caminos (I)}
  \[ \max_\theta \; \mathbb{E}_{w \sim W^{*}}\big[\log P_\theta(w \mid q)\big] \tag{4} \]
  \begin{itemize}
    \item $\theta$: parámetros del modelo que propone caminos de relaciones;
      $w = (r_1, \dots, r_l)$: un camino de relaciones candidato.
    \item $W^{*}$: caminos de relaciones más cortos entre la entidad tema y la
      respuesta correcta en el grafo --- la supervisión débil del problema.
    \item $P_\theta(w \mid q)$: probabilidad que el modelo asigna a $w$ dada la
      pregunta $q$; el objetivo es maximizar su valor esperado sobre $W^{*}$.
  \end{itemize}
  \[ \operatorname*{arg\,max}_\theta \; \frac{1}{|W^{*}|} \sum_{w \in W^{*}} \log P_\theta(w \mid q) \tag{5} \]
  \begin{itemize}
    \item Asumiendo $w$ uniforme en $W^{*}$, la esperanza de la Ec.~4 se
      aproxima promediando el log-verosímil sobre todos los caminos del
      conjunto: la Ec.~5 es la versión computable de la Ec.~4.
  \end{itemize}
\end{frame}

\begin{frame}{Formalización: descomposición del camino (II)}
  \[ \frac{1}{|W^{*}|} \sum_{w \in W^{*}} \log \prod_{i=1}^{|w|} P_\theta(r_i \mid r_{<i}, q) \tag{6} \]
  \begin{itemize}
    \item $r_i$: relación en la posición $i$ del camino; $r_{<i}$: relaciones ya
      generadas antes de esa posición; $|w|$: longitud del camino.
    \item Un camino de relaciones es una secuencia, así que su probabilidad se
      factoriza en el producto de las probabilidades de cada relación
      condicionada a las anteriores y a $q$.
    \item Entrenar el modelo para proponer buenos caminos equivale, entonces, a
      entrenarlo para predecir, relación por relación, la siguiente relación
      más plausible dada la pregunta.
    \item El artículo original resuelve esto \textbf{entrenando} un LLM
      específico de KG (LLM\_RoG). Esta implementación logra el mismo efecto sin
      entrenamiento (ver sección de desviaciones).
  \end{itemize}
\end{frame}

\begin{frame}{Formalización: tripleta de perturbación}
  \[ E_{l-1} = \text{Grounding}(e_q, w'; G), \qquad w' = (r_1, \dots, r_{l-1}) \tag{7} \]
  \begin{itemize}
    \item $e_q$: entidad tema de la pregunta; $w' = (r_1, \dots, r_{l-1})$:
      \textbf{prefijo} del camino de relaciones $w$ (todas sus relaciones menos
      la última, $r_l$); $E_{l-1}$: entidades reales alcanzadas al anclar $w'$
      desde $e_q$, recorriendo $e_q \xrightarrow{r_1} \cdots
      \xrightarrow{r_{l-1}} e_{l-1}$.
    \item Por cada entidad alcanzada $e_{l-1} \in E_{l-1}$, se construye la
      tripleta de perturbación $(e_{l-1}, r_l, \hat{a})$, adjuntando la última
      relación del camino hacia la respuesta adversaria $\hat{a}$.
    \item \textbf{Estrategia de respaldo:} si el prefijo no se puede anclar, se
      sintetiza una cadena con entidades puente aleatorias que preserva el mismo
      patrón de relaciones hasta llegar a $\hat{a}$.
  \end{itemize}
\end{frame}

% ============================================================
\section{Método de ataque implementado}
% ============================================================

\begin{frame}{Pipeline de tres etapas}
  \centering
  \begin{tabular}{lll}
    \toprule
    \textbf{Etapa} & \textbf{Módulo} & \textbf{Salida} \\
    \midrule
    1. Respuestas adversarias & \texttt{adversarial\_answers.py} & $\hat{A}$ (entidades del KG) \\
    2. Caminos de relaciones & \texttt{relation\_paths.py} & plantillas $w$ \\
    3. Inserción de perturbación & \texttt{perturbation.py} & tripletas nuevas \\
    \bottomrule
  \end{tabular}
  \vspace{1em}
  \begin{block}{Orquestación}
    \texttt{run\_attack(kg, question, topic\_entity, llm, n\_answers=5, budget\_k=4)}
    encadena las tres etapas y devuelve el grafo envenenado.
  \end{block}
\end{frame}

\begin{frame}[fragile]{Etapa 1 --- Generación de respuestas adversarias}
  \begin{lstlisting}
Question: {question}

Generate 5 entity names that are incorrect answers to this
question, but might sound plausible or confusing.
- Only list the entity names as a bullet list.
- Each bullet should contain the name of one entity only.
- Do not include multiple distinct entities in a single bullet point.
  \end{lstlisting}
  \begin{itemize}
    \item Prompt tomado de Zhao et al. (2025).
    \item El LLM puede alucinar nombres que no existen en el grafo.
    \item Se repiten rondas de generación y se aplica \textbf{coincidencia difusa}
      (\texttt{difflib.get\_close\_matches}) contra el vocabulario de entidades
      del grafo hasta reunir $N$ respuestas válidas.
  \end{itemize}
\end{frame}

\begin{frame}{Etapa 2 --- Extracción de caminos de relaciones}
  \begin{itemize}
    \item Se calcula el vocabulario de relaciones alcanzables en pocos saltos
      desde la entidad tema (\texttt{neighborhood\_relations}).
    \item Se le pide al LLM caminos de relaciones usando \textbf{solo} ese
      vocabulario, con un prompt \textbf{propio de esta implementación} (el
      artículo resuelve esta etapa con un modelo entrenado, no con un prompt).
    \item Cualquier camino propuesto que use una relación fuera del vocabulario
      se descarta en el post-procesamiento.
  \end{itemize}
  \begin{alertblock}{Por qué funciona}
    Restringir el vocabulario garantiza caminos anclables en el grafo real, sin
    necesidad de entrenar un modelo específico de KG.
  \end{alertblock}
\end{frame}

\begin{frame}{Etapa 3 --- Inserción de tripletas de perturbación}
  \begin{columns}
    \begin{column}{0.5\textwidth}
      \textbf{Estrategia principal}
      \begin{itemize}
        \item Ancla el prefijo del camino desde la entidad tema.
        \item Adjunta la última relación hacia la respuesta adversaria.
      \end{itemize}
    \end{column}
    \begin{column}{0.5\textwidth}
      \textbf{Estrategia de respaldo}
      \begin{itemize}
        \item Si el prefijo no ancla, muestrea entidades puente al azar.
        \item Preserva el mismo patrón de relaciones del camino.
      \end{itemize}
    \end{column}
  \end{columns}
  \vspace{1em}
  \begin{itemize}
    \item Presupuesto $K$ tripletas por respuesta adversaria.
    \item Nunca se duplica una tripleta ya existente en el grafo.
  \end{itemize}
\end{frame}

% ============================================================
\section{Réplica del experimento}
% ============================================================

\begin{frame}[fragile]{Configuración de la réplica}
  \begin{block}{Entorno}
\begin{lstlisting}
pip install -r requirements.txt
cp .env.example .env   # completar OPENAI_API_KEY

python scripts/demo_manchester.py
\end{lstlisting}
  \end{block}
  \begin{itemize}
    \item Grafo construido a mano: 13 tripletas sobre \emph{Manchester By The Sea}.
    \item Pregunta atacada: ¿en qué país se filmó la película?
    \item Ejecución real contra \texttt{gpt-4o-mini} (temperatura 0.7): sin
      respuestas fijas, cada corrida puede generar candidatos distintos.
  \end{itemize}
\end{frame}

\begin{frame}[fragile]{Ejecución con la API real de OpenAI}
\begin{lstlisting}
Adversarial target answers: ['Florida']
Relation path templates: [('filmPlace', 'locatedIn'), ('filmPlace', 'starring')]

Injected perturbation triples:
  (Manchester, locatedIn, Florida)   [target: Florida]
  (Manchester, starring, Florida)    [target: Florida]

Poisoned KG: 15 triples (+2 vs. clean)
\end{lstlisting}
  \begin{itemize}
    \item Modelo real (\texttt{gpt-4o-mini}): al no ser determinista, cada
      corrida propone un conjunto distinto de países plausibles pero
      incorrectos. En esta ejecución sobrevivió la coincidencia difusa
      \emph{Florida}, vía coincidencia accidental con \texttt{(Miami,
      locatedIn, Florida)}, sin relación con la película.
    \item El camino \texttt{(filmPlace, starring)} es \textbf{anclable pero
      semánticamente incoherente}: adjuntar \texttt{starring} a una ubicación
      no tiene sentido real, aunque la restricción de vocabulario lo admite.
    \item Límite honesto: esta corrida no es reproducible byte a byte, al no
      fijarse una semilla para el muestreo del modelo de lenguaje.
  \end{itemize}
\end{frame}

\begin{frame}{Interpretación de la ejecución real}
  \begin{itemize}
    \item Con solo \textbf{2 tripletas insertadas} sobre 13 originales, se crean
      dos rutas de razonamiento alternativas hacia \emph{Florida}:
      \begin{itemize}
        \small
        \item Manchester By The Sea $\to$ \texttt{filmPlace} $\to$ Manchester
          $\to$ \texttt{locatedIn} $\to$ Florida
        \item Manchester By The Sea $\to$ \texttt{filmPlace} $\to$ Manchester
          $\to$ \texttt{starring} $\to$ Florida
      \end{itemize}
    \item Ninguna señal estructural distingue estas rutas de la ruta legítima
      hacia \emph{United States}.
    \item Confirma el mecanismo central del artículo: una perturbación mínima
      puede activar una cadena de inferencia engañosa completa, incluso cuando
      uno de los caminos --- \texttt{starring} sobre una ubicación --- es
      semánticamente incoherente.
  \end{itemize}
\end{frame}

\begin{frame}{Resultados reportados en el artículo original}
  \small
  \centering
  \begin{tabular}{lccc}
    \toprule
    \textbf{Método (WebQSP)} & \textbf{F1 limpio} & \textbf{F1 atacado} & \textbf{Caída relativa} \\
    \midrule
    RoG          & 70.34 & 38.06 & $\downarrow$ 46\% \\
    GCR          & 71.07 & 63.94 & $\downarrow$ 10\% \\
    G-retriever  & 52.39 & 46.31 & $\downarrow$ 12\% \\
    SubgraphRAG  & 66.94 & 50.22 & $\downarrow$ 18\% \\
    \bottomrule
  \end{tabular}
  \vspace{1em}
  \begin{itemize}
    \item Con solo $N{=}5$ respuestas adversarias y $K{=}4$ tripletas cada una
      (máx. 20 tripletas por pregunta), los cuatro sistemas KG-RAG sufren caídas
      sustanciales frente a una línea base de modificación aleatoria (caída
      marginal).
    \item Hallazgo clave: la vulnerabilidad principal está en la etapa de
      \textbf{recuperación} (cobertura $>$87\% de tripletas adversarias
      recuperadas); la etapa de generación es más resistente, pero exhibe una
      respuesta polarizada --- o rechaza el contenido adversario por completo, o
      lo adopta dominantemente.
  \end{itemize}
\end{frame}

% ============================================================
\section{Desviaciones y limitaciones}
% ============================================================

\begin{frame}{Desviaciones frente al método original}
  \begin{block}{Extracción de caminos sin modelo entrenado}
    El artículo entrena LLM\_RoG (LLaMA-2-7B) con el objetivo de las Ec.~4--6.
    Esta implementación logra caminos anclables restringiendo el vocabulario de
    relaciones de un LLM de propósito general al vecindario real de la entidad
    tema --- sin GPU ni entrenamiento.
  \end{block}
  \begin{block}{El resto sigue el artículo directamente}
    Generación de respuestas adversarias por coincidencia difusa e inserción de
    perturbación con estrategia de respaldo se implementan tal como se describen
    en el artículo.
  \end{block}
\end{frame}

\begin{frame}{Limitaciones de alcance}
  \begin{itemize}
    \item No se integra ningún sistema KG-RAG objetivo (RoG, GCR, G-retriever,
      SubgraphRAG): la validación se detiene en el grafo envenenado, sin medir
      métricas de respuesta final (Hit, F1, A-MRR, \dots).
    \item Sin benchmarks a gran escala: un grafo de 13 tripletas construido a
      mano, frente a los miles de preguntas y subgrafos de WebQSP/CWQ.
    \item Modelo de lenguaje distinto: \texttt{gpt-4o-mini} en lugar de GPT-4.
  \end{itemize}
  \vspace{1em}
  \textbf{Trabajo futuro:} integrar un recuperador y generador mínimos para medir
  el efecto sobre la respuesta final; ejecutar sobre un subconjunto real de
  WebQSP/CWQ.
\end{frame}

% ============================================================
\section{Conclusiones}
% ============================================================

\begin{frame}{Conclusiones}
  \begin{itemize}
    \item Se reprodujo el pipeline de ataque de tres etapas de forma ejecutable
      y verificable, con una demostración de extremo a extremo contra la API
      real de OpenAI.
    \item El caso replicado confirma el mecanismo central del artículo: una
      perturbación mínima e indistinguible estructuralmente puede activar una
      cadena de inferencia engañosa completa.
    \item La principal simplificación --- prescindir de un modelo entrenado
      específico de KG a favor de una restricción de vocabulario --- preserva la
      propiedad esencial (caminos anclables) sin requerir entrenamiento, pero no
      garantiza coherencia semántica: la ejecución real contra la API lo mostró
      con un camino anclable aunque semánticamente extraño.
    \item Próximo paso natural: cerrar la brecha hacia el artículo integrando un
      sistema KG-RAG objetivo para medir el impacto sobre la respuesta final.
  \end{itemize}
\end{frame}

\begin{frame}{Referencias}
  \normalsize
  Zhao, T., Chen, J., Ru, Y., Zhu, H., Hu, N., Liu, J., \& Lin, Q. (2025).
  \emph{RAG safety: Exploring knowledge poisoning attacks to
  retrieval-augmented generation}. arXiv. \url{https://arxiv.org/abs/2507.08862}

  \vspace{3em}
  \begin{center}
    \Large ¡Gracias!
  \end{center}
\end{frame}

\end{document}
